\usepackage[top=2.5cm, bottom=3cm, left=3cm, right=3cm]{geometry}
\usepackage[spanish, main=english]{babel} %-----tildes  [spanish,activeacute,es-tabla]
\usepackage[pdftex]{graphicx}
\usepackage{amsmath, amssymb}%-----estructuras matemáticas
%
\usepackage{textcase}
\usepackage{caption}
\DeclareCaptionTextFormat{up}{\MakeTextUppercase{#1}}
\captionsetup[table]{labelformat={default}, font=footnotesize, justification=centering, labelsep=newline, name={TABLA }, textformat=up} %-----formato del caption table
\renewcommand{\thetable}{\Roman{table}}
\captionsetup[figure]{labelformat={default}, font=footnotesize, labelsep=period,  name={Fig. }, justification=centering,singlelinecheck=false} %-----formato del caption figure
\renewcommand{\thefigure}{\arabic{figure}}
%
\usepackage{hyperref}
   \hypersetup{colorlinks, linkcolor=blue, urlcolor=blue, citecolor={red}}
\usepackage{pstricks-add}
\usepackage{longtable}
\usepackage{pst-all}
\usepackage{verbatim}%-----para mostrar comnados 
\usepackage{hhline, colortbl}
\usepackage{color}
\usepackage{array}
\usepackage{enumitem}



\usepackage[numbers]{natbib}

\setlength{\parindent}{1.25cm}%sangria
\renewcommand{\baselinestretch}{1.5}%interlineado
\setlength{\parskip}{0.1em}%espacio entre parrafos


\renewcommand{\labelitemi}{$\bullet $}
\setlist[itemize]{leftmargin=*, parsep=-5px}
\renewcommand\descriptionlabel[1]{$\bullet$\hspace{0.5em} \textbf{#1}}
\setlist[description]{parsep=-5px}

\usepackage{tocloft}
\renewcommand{\cftfigfont}{Fig. }
\renewcommand{\cfttabfont}{Tabla }
\tocloftpagestyle{empty}

\usepackage{titletoc}
\titlecontents{section}[0.1em]{\smallskip}{\enspace\thecontentslabel.\hspace{1em}} {\hspace*{}} {\titlerule*[1pc]{.}\contentspage}
\titlecontents{subsection}[0.1em]{\smallskip}{\hspace{0.5cm}\enspace\thecontentslabel.\hspace{1em}} {\hspace*{}} {\titlerule*[1pc]{.}\contentspage}
\titlecontents{subsubsection}[0.1em]{\smallskip}{\hspace{1cm}\enspace\thecontentslabel)\hspace{1em}} {\hspace*{}} {\titlerule*[1pc]{.}\contentspage}
\titlecontents{paragraph}[0.1em]{\smallskip}{\hspace{1.5cm}\enspace\thecontentslabel)\hspace{1em}} {\hspace*{}} {\titlerule*[1pc]{.}\contentspage}



\setcounter{tocdepth}{5}
\setcounter{secnumdepth}{4}%define el contador para los 4 niveles de titulo


% \renewcommand{\thesection}{\Roman{section}}
% \renewcommand{\thesubsection}{\Alph{subsection}}   
% \renewcommand{\thesubsubsection}{\arabic{subsubsection}}
% \renewcommand{\theparagraph}{\alph{paragraph}}



\usepackage{titlesec}
\titleformat{\section}{\centering\normalfont\normalsize\MakeUppercase}{\thesection.}{0.2em}{}
\titleformat{\subsection}{\normalfont\normalsize\it}{\thesubsection.}{0.2em}{}
\titleformat{\subsubsection}[runin]{\itshape\normalsize\it}{\hspace{0.5cm}\thesubsubsection)}{0.2em}{}
\titleformat{\paragraph}[runin]{\normalfont\normalsize\it}{\hspace{1cm}\theparagraph)}{0.2em}{}
%\titleformat{\subparagraph}[runin]{\itshape\normalsize\bfseries}{\thesubparagraph}{1em}{}


% \usepackage{fancyhdr}

% \definecolor{verdeUdeA}{rgb}{0.325,0.506,0.208} %rgb=83,129,53

% \setlength{\headheight}{25pt}
% \renewcommand{\headrule}{{\color{verdeUdeA}\rule{\textwidth}{1pt}}}
% \renewcommand{\headrulewidth}{2pt}
% \lhead{\footnotesize\MakeUppercase{\smalltitulo }}\chead{}\rhead{{\color{white}Rufino}\thepage} %borrar los encabezados
% \cfoot{}

% --- Configuración de Encabezados Dinámicos ---
% \usepackage{fancyhdr}

% \definecolor{verdeUdeA}{rgb}{0.325,0.506,0.208} %rgb=83,129,53

% \pagestyle{fancy}
% % \fancyhf{} % Limpia todos los encabezados y pies de página

% % Definimos que \leftmark solo guarde el nombre del capítulo (sin la palabra "Chapter X")
% \renewcommand{\chaptermark}[1]{\markboth{#1}{}} 

% % Configuración del encabezado
% \lhead{\footnotesize\MakeUppercase{\leftmark}} % <--- Aquí cambiamos \smalltitulo por \leftmark
% \chead{}
% \rhead{\thepage} % El número de página a la derecha

% % Línea verde de la UdeA
% \renewcommand{\headrule}{{\color{verdeUdeA}\rule{\textwidth}{1pt}}}
% \renewcommand{\headrulewidth}{2pt}

% % Aseguramos que el pie de página esté vacío
% \cfoot{}

% --- Configuración de Encabezados Dinámicos e Inteligentes ---
\usepackage{fancyhdr}
\usepackage{etoolbox} % Necesario para la lógica IF

\definecolor{verdeUdeA}{rgb}{0.325,0.506,0.208} 

\pagestyle{fancy}
\fancyhf{} % Limpia encabezados y pies anteriores

% 1. Definimos cómo se guarda la marca del capítulo
% Esto hace que \leftmark contenga solo el título del capítulo
\renewcommand{\chaptermark}[1]{\markboth{#1}{}} 

% 2. Lógica del Encabezado Izquierdo (\lhead)
% IF: Si \leftmark está vacío (antes del primer capítulo), muestra \smalltitulo
% ELSE: Si ya hay un capítulo, muestra \leftmark
\lhead{\footnotesize\MakeUppercase{%
  \expandafter\ifstrempty\expandafter{\leftmark}{%
    \smalltitulo%
  }{%
    \leftmark%
  }%
}}

\chead{}
\rhead{\thepage} % Número de página a la derecha

% 3. Línea verde de la UdeA
\renewcommand{\headrule}{{\color{verdeUdeA}\rule{\textwidth}{1pt}}}
\renewcommand{\headrulewidth}{2pt}
\cfoot{}

